\subsection{Convergence of the hierarchical construction in the dense graph limit}
\label{sec:convergence_dense}

The hierarchical construction of $\chi_\mathrm{eff}$ in Appendix~\ref{subsec:chi_eff}
is defined for a fixed finite graph.
A natural question is whether this construction commutes with the limit in which the
graph becomes dense — that is, whether the effective field and the induced spectral
distances converge to well-defined continuous objects as the number of vertices $N$
grows and the neighborhood radius $\ell_0$ is scaled appropriately.

A complete proof of such convergence lies outside the scope of this work.
However, the construction is compatible with the known convergence results for graph
Laplacians in the following sense.
The combinatorial neighborhood $\mathcal{N}_i$ used to define $\bar\chi$ has radius
$\ell_0$ in combinatorial distance, which corresponds to a ball of radius
$O(\ell_0 / N^{1/d})$ in the ambient metric as the graph becomes dense.
In this regime, $\bar\chi_i \to \chi(x_i)$ pointwise, where $\chi$ is the continuous
field from which the discrete configuration is sampled.
The weighted Laplacian constructed from these $\bar\chi$ values then converges to
the Laplace--Beltrami operator weighted by the local relational stiffness, by the
results of Belkin and Niyogi~\cite{Belkin2003} and Hein, Audibert, and von
Luxburg~\cite{Hein2007}, provided $\ell_0 = \ell_0(N)$ is chosen to satisfy the
standard bandwidth conditions.
Under these conditions, the effective field $\chi_\mathrm{eff}$ converges to a
smooth coarse-grained version of $\chi$ on the underlying manifold, and the spectral
distances $d_S(i,j)$ approximate the continuous diffusion distances.

We therefore conclude that the finite-graph construction is consistent with the
continuous limit under standard regularity assumptions, and that no structural
obstruction arises from the hierarchical dependency chain in this limit.
